\subsection{A short history of AI}
\label{sec:ai-history}
\slidesrc{Thinking machines: 13--67}

The slides build a timeline one event at a time (slides 16--67). The complete timeline is
redrawn in Figure~\ref{fig:ai-timeline}; the sub-sections below explain each step.

\begin{center}
\begin{tikzpicture}[font=\small,
  ev/.style={anchor=north west, text width=10.6cm, align=left, inner sep=2pt},
  yr/.style={anchor=north east, font=\small\bfseries, inner sep=2pt}]
  \coordinate (e0) at (0.3,0);
  \foreach \i/\year/\txt/\c in {
    1/1950/{Alan Turing proposes the ``Turing test'' to determine a machine's human-like ability}/nbblue,
    2/1956/{The term ``Artificial Intelligence'' is coined at the Dartmouth workshop}/nbblue,
    3/1966--72/{Stanford Research Institute builds Shakey, the first robot that reasons about its own actions}/nbgreen,
    4/1970--72/{Colmerauer designs Prolog; first Prolog system by Colmerauer and Roussel}/nbgreen,
    5/1974--80/{First ``AI winter'': the Lighthill report (1973) cuts funding}/nbred,
    6/1980/{Symbolics Lisp machines are commercialised: AI renaissance (expert systems)}/nbgreen,
    7/1987--93/{Second ``AI winter'': collapse of the Lisp machine market}/nbred,
    8/1997/{IBM Deep Blue defeats world chess champion Garry Kasparov}/nbgreen,
    9/2008/{Google speech-to-search iPhone app}/nborange,
    10/2011/{IBM Watson defeats Ken Jennings at the quiz show Jeopardy!}/nborange,
    11/2011--12/{Apple introduces Siri, an intelligent personal assistant}/nborange,
    12/2012/{Google Brain learns to identify cats on YouTube (75\% accuracy) without being told what a cat is}/nborange,
    13/2013/{China's Tianhe-2 doubles the world's top supercomputing speed (33.86 petaflops)}/nbgray,
    14/2014/{Facebook develops DeepFace, face recognition with near-human accuracy}/nborange,
    15/2015/{Google open-sources the TensorFlow machine learning library}/nborange,
    16/2016/{Google DeepMind's AlphaGo defeats Lee Sedol at Go}/nborange,
    17/2016/{China's Sunway TaihuLight triples the top supercomputing speed (93 petaflops)}/nbgray,
    18/2016/{Uber pilots a self-driving car programme in Pittsburgh, PA}/nborange,
    19/2019--now/{AlphaStar, AlphaFold, generative AI, agentic AI (Section~\ref{sec:recent-achievements})}/nbpurple}
  {
    \pgfmathtruncatemacro{\prev}{\i-1}
    \node[ev] (e\i) at ([yshift=-3pt]e\prev.south west) {\txt};
    \node[yr, text=\c] at ([xshift=-0.55cm]e\i.north west) {\year};
    \fill[\c] ([xshift=-0.3cm, yshift=-7pt]e\i.north west) circle (3.2pt);
  }
  \begin{scope}[on background layer]
    \draw[very thick, nbblue!50] (0,0) -- (0,0 |- e19.south);
  \end{scope}
\end{tikzpicture}
\captionof{figure}{The AI timeline of the slides (slide 67, extended). Colours:
\textcolor{nbblue}{foundations}, \textcolor{nbgreen}{symbolic AI successes},
\textcolor{nbred}{AI winters}, \textcolor{nborange}{data-driven / neural successes},
\textcolor{nbgray}{computing power}, \textcolor{nbpurple}{today}. Notice the shift: until the
1990s the big results are symbolic (reasoning, planning, chess search); after 2008 they are
mostly learned from data.}
\label{fig:ai-timeline}
\end{center}

\subsubsection{Before computers}
\slidesrc{Thinking machines: 13--15}

The slides show three portraits, without text. Here is why each person is on the
``before'' list (the explanations are added):
\begin{itemize}[nosep]
  \item \textbf{Aristotle} (384--322 BC) invented the first system of formal logic, the
        \emph{syllogism}: ``All men are mortal; Socrates is a man; therefore Socrates is
        mortal''. The very same example will come back in first-order logic
        (Whiteboard~\ref{wb:modus-ponens}) and in Prolog.
  \item \textbf{Gottfried Wilhelm von Leibniz} (1646--1716) dreamed of a universal symbolic
        language (\emph{characteristica universalis}) and a ``calculus of reasoning''
        (\emph{calculus ratiocinator}), so that disputes could be settled by saying
        ``let us calculate!''. He also studied binary arithmetic and built a mechanical
        calculator. This is the dream of symbolic AI: reasoning as computation.
  \item \textbf{Mary Shelley} (1797--1851) wrote \emph{Frankenstein} (1818), the story of an
        artificial being created by a scientist, and of the consequences of creating it:
        the first great story about the risks of artificial creatures
        (see Section~\ref{sec:risks-and-regulation}).
\end{itemize}

\subsubsection{1950: the Turing test}
\slidesrc{Thinking machines: 16--23}

In 1950 Alan Turing published the paper \emph{Computing Machinery and Intelligence} in the
philosophy journal \emph{Mind}. Instead of answering the vague question ``Can machines
think?'', he replaced it with a game that can actually be played, the
\textbf{imitation game}, now called the \textbf{Turing test}.

\begin{nbdef}{Turing test}
A human interrogator chats (in writing) with two hidden partners: a human and an AI system.
By asking questions, the interrogator must decide which one is the machine. If the
interrogator cannot reliably tell them apart, the machine passes the test.
\end{nbdef}

\begin{center}
\begin{tikzpicture}[>=Stealth, font=\small,
  who/.style={draw, rounded corners, thick, minimum width=2.4cm, minimum height=0.8cm, align=center}]
  \node[who, draw=nbblue, fill=nbblue!8] (int) at (0,0) {Human\\interrogator};
  \draw[ultra thick, nbgray] (2.2,-1.6) -- (2.2,1.6) node[above, font=\scriptsize]{wall};
  \node[who, draw=nbgreen, fill=nbgreen!8] (h) at (5,0.85) {Human (A or B?)};
  \node[who, draw=nbred, fill=nbred!8] (m) at (5,-0.85) {AI system (A or B?)};
  \draw[<->, thick] (int.east) -- (2.2,0) node[right, font=\Large]{?};
  \draw[<->, thick] (2.5,0) -- (h.west);
  \draw[<->, thick] (2.5,0) -- (m.west);
\end{tikzpicture}
\captionof{figure}{The imitation game (slide 22): the interrogator only sees written answers and
must guess which partner is the machine.}
\end{center}

\begin{nbintuition}[Behaviour instead of definitions]
Turing's trick is to judge intelligence by \emph{observable behaviour}, not by what happens
inside. The agent definition of AI (Section~\ref{sec:ai-definition}) has the same spirit.
The slides add a comic (``Turing test 2208'') where the human fails: ``the machine realised
immediately that it was me being the dumb human''. Today large language models make the
classic Turing test much less decisive than it looked in 1950.
\end{nbintuition}

\paragraph{An open parenthesis (slides 19--21).} What does it exactly mean to ``think'', to
``perceive'', to be ``conscious of being alive''? The slides link to the neuroscientist
Anil Seth (TED talk \emph{Your brain hallucinates your conscious reality}): according to him,
perception is a \emph{controlled hallucination}, the brain's best guess about the causes of
its sensory signals, constantly corrected by the senses. If we cannot define these words
precisely for humans, judging them in machines is even harder. (Closed parenthesis.)

\subsubsection{1956: the Dartmouth workshop and the founding fathers}
\slidesrc{Thinking machines: 24--33}

In summer 1956 a workshop at Dartmouth College (USA) coined the term
\textbf{``Artificial Intelligence''}. Its proposal contained the founding conjecture of the
field:

\begin{center}
\begin{minipage}{0.88\linewidth}
\itshape ``Every aspect of learning or any other feature of intelligence can in principle be
so precisely described that a machine can be made to simulate it.''
\end{minipage}
\end{center}

The slides list the participants (the ``founding fathers of AI'') and, step by step, what
each of them is remembered for:

\begin{center}
\small
\begin{tabularx}{\linewidth}{@{}l X@{}}
\toprule
\textbf{Person} & \textbf{Known for} \\
\midrule
John McCarthy & the \emph{Advice Taker} (a program reasoning with logical sentences, which
                inspired Prolog), the Lisp language, ALGOL \\
Marvin Minsky & artificial neural networks, the \emph{society of mind} (the mind as a
                society of simple interacting agents: an ancestor of multi-agent systems!) \\
Claude Shannon & information theory, cryptography, natural language processing \\
Ray Solomonoff & machine learning, prediction, Bayesian frameworks \\
Allen Newell \& Herbert Simon & the \emph{Logic Theory Machine} (proved theorems of
                \emph{Principia Mathematica}) and the \emph{General Problem Solver} \\
Arthur Samuel & search, alpha-beta pruning, machine learning (a checkers program that
                learned to play better than its author) \\
\midrule
Oliver Selfridge & Pandemonium theory (pattern recognition by many competing ``demons'') \\
Nathaniel Rochester & IBM, designer of the IBM 701 computer \\
Trenchard More & natural deduction (the proof system of Section~\ref{sec:natural-deduction-prop}) \\
\bottomrule
\end{tabularx}
\captionof{table}{The founding fathers of AI and their contributions (slides 26--33).}
\end{center}

\begin{nbintuition}[Symbolic and sub-symbolic, from day one]
Already at Dartmouth both souls of AI were present: logic and symbols (McCarthy, Newell and
Simon, More) and learning and neural networks (Minsky, Solomonoff, Samuel). The rest of the
history is a pendulum between the two.
\end{nbintuition}

\subsubsection{1966--1972: Shakey the robot}
\slidesrc{Thinking machines: 34--36}

\textbf{Shakey} was the first general-purpose mobile robot able to \emph{reason about its own
actions}. Other robots had to be instructed on each individual step; Shakey could analyse a
command and break it down into basic chunks by itself. The project combined robotics,
computer vision and natural language processing, and was the first to meld \emph{logical
reasoning} and \emph{physical action}. It was developed at the AI Center of the Stanford
Research Institute (now SRI International) from about 1966 to 1972 (project manager Charles
Rosen), programmed mainly in LISP, and its planning component was the \textbf{STRIPS}
planner. As the first logical, \textbf{goal-based agent}, Shakey lived in a limited world
(a few rooms with boxes).

\begin{nbexample}{what ``breaking a command into chunks'' means}
\notinslides
Command: ``push the box from room A to room B''. A planner like STRIPS knows actions with
preconditions and effects, e.g.\ \code{goto(Place)}, \code{push(Box, From, To)}, and finds
the sequence: \code{goto(room\_A)} $\to$ \code{goto(box)} $\to$
\code{push(box, room\_A, door)} $\to$ \code{push(box, door, room\_B)}. Your navigation app
and a warehouse robot do the same kind of planning today. Shakey is a deliberative agent
in the sense of Section~\ref{sec:agent-architectures}.
\end{nbexample}

\subsubsection{1970--1972: Prolog}
\slidesrc{Thinking machines: 37}

In 1970 Alain Colmerauer designed the \textbf{Prolog} language (from the French
\emph{programmation en logique}); the first Prolog system was implemented by Colmerauer and Philippe Roussel in 1972
in Marseille. Prolog is the language of Part~IV of these notes: a program is a set of
logical facts and rules, and running it means proving a query.

\subsubsection{The AI winters and the Lisp machines}
\slidesrc{Thinking machines: 38--43}

Early AI promised a lot. When the promises were not kept, funding collapsed: these periods
are called \textbf{AI winters}.

\begin{itemize}
  \item \textbf{First AI winter (1974--1980).} The \textbf{Lighthill report} (\emph{Artificial
        Intelligence: A General Survey}, Sir James Lighthill, 1973, for the British Science
        Research Council) gave a very pessimistic evaluation: ``in no part of the field have
        discoveries made so far produced the major impact that was then promised''. It formed
        the basis for the British government's decision to end support for AI research in
        all but three universities (Edinburgh, Sussex, Essex). Its main technical criticism:
        AI had failed to address the \textbf{combinatorial explosion} that appears when
        problems are taken from toy worlds to real-world domains.
  \item \textbf{AI renaissance (1980).} Symbolics commercialises \textbf{Lisp machines}:
        computers with hardware support for Lisp operations (data type testing, CDR coding,
        incremental garbage collection), running large Lisp programs very efficiently, with
        operating systems written entirely in Lisp. They powered the boom of
        \emph{expert systems} (rule-based programs encoding the knowledge of human experts).
  \item \textbf{Second AI winter (1987--1993).} With the microcomputer revolution, cheap
        desktop PCs soon ran Lisp programs even faster than Lisp machines, without special
        hardware. The Lisp machine market collapsed: most makers went out of business by the
        early 1990s, and AI lost credibility again.
\end{itemize}

\begin{nbexample}{combinatorial explosion in numbers}
\notinslides
Chess has about 35 legal moves per position. Looking $d$ moves ahead means examining about
$35^d$ positions:
\[
35^2 = 1\,225, \qquad 35^4 = 1\,500\,625, \qquad 35^6 \approx 1.8 \times 10^{9}, \qquad
35^{10} \approx 2.8 \times 10^{15}.
\]
Each extra level multiplies the work by 35. An algorithm that works on a toy problem can
become hopeless on a slightly bigger one: this is what the Lighthill report meant. Deep Blue
(1997) beat it on chess with special hardware plus pruning (alpha-beta, Samuel's legacy).
\end{nbexample}

\subsubsection{1997--2016: from chess to Go, the data-driven era}
\slidesrc{Thinking machines: 44--67}

\begin{itemize}[nosep]
  \item \textbf{1997}: IBM \textbf{Deep Blue} defeats the world chess champion Garry
        Kasparov (symbolic AI: search plus hand-written evaluation).
  \item \textbf{2008}: Google releases a \textbf{speech-to-search} iPhone app: you speak, the
        recording is sent to Google's servers, and the results come back in seconds.
  \item \textbf{2011}: IBM \textbf{Watson} defeats the champion Ken Jennings at the quiz show
        \emph{Jeopardy!} (natural language question answering).
  \item \textbf{2011/2012}: Apple introduces \textbf{Siri}, an intelligent personal assistant.
  \item \textbf{2012}: after seeing 10 million images from YouTube videos in three days, the
        16{,}000-computer network of \textbf{Google Brain} (one billion connections) began to
        recognise cats, even though it had never been taught what a cat looked like
        (unsupervised learning, 75\% accuracy).
  \item \textbf{2013}: China's Tianhe-2 doubles the world's top supercomputing speed (33.86
        petaflops); \textbf{2016}: Sunway TaihuLight triples it (93 petaflops). Raw computing
        power is one of the reasons for the success of deep learning.
  \item \textbf{2014}: Facebook's \textbf{DeepFace} identifies human faces in images: a
        nine-layer neural network with over 120 million connection weights, trained on four
        million images uploaded by Facebook users, said to be 97\% accurate (vs.\ 85\% for
        the FBI's Next Generation Identification system). It uses a 3D model to virtually
        rotate faces so that they look at the camera.
  \item \textbf{2015}: Google open-sources \textbf{TensorFlow}, a library for numerical
        computation with data flow graphs, now widely used for machine learning.
  \item \textbf{2016}: Google DeepMind's \textbf{AlphaGo} defeats Lee Sedol at Go, a game
        whose search space is far too large for Deep Blue-style brute force.
  \item \textbf{2016}: Uber pilots a \textbf{self-driving car} programme in Pittsburgh, PA.
\end{itemize}

\begin{nbwarning}[Two kinds of ``winning'']
Deep Blue and AlphaGo both beat a world champion, but in very different ways. Deep Blue
\emph{searched} millions of positions per second with rules written by humans (System~2
style, symbolic). AlphaGo \emph{learned} an intuition for good moves from millions of games
with neural networks, combined with a search (System~1 + System~2). Do not file both under
``AI wins at games'' without noting the difference: it is the theme of
Section~\ref{sec:symbolic-vs-subsymbolic}.
\end{nbwarning}
